\section{A growing-modulus arithmetic application: coprime circles}
\label{sec:coprimecircles}
This application uses exact finite Fourier analysis, with no
prime number theorem, fixed-modulus assumption or prime-tuple input.
Let \(X_n\) have vertex set \(\mathbb Z/n\mathbb Z\), with
distinct vertices adjacent when their difference is coprime to \(n\).
These are the unitary Cayley graphs.
Their Hamiltonicity and stronger path properties are known;
see Klotz--Sander~\cite{KS} and Ba\v{s}i\'c~\cite{Basic}.
The conclusion here concerns the full exponential count.

\begin{theorem}[Coprime-circle entropy at every modulus]
\label{thm:coprimecircles}
As \(n\to\infty\) through all integers,
\[
 \log h(X_n)
   =n(\log\varphi(n)-1)
       +O\left(n\varphi(n)^{-1/8}\log(2n)\right).
\]
In particular \(h(X_n)=(\varphi(n)/e)^n\exp(o(n))\).
The constant is absolute.
For odd \(n\), the same logarithmic estimate holds for the
number of spanning paths joining any prescribed distinct endpoints.
For even \(n\), it holds for any prescribed opposite-parity endpoints.
\end{theorem}
\begin{proof}
The degree is \(d=\varphi(n)\), and \(P=A_{X_n}/d\) is symmetric
stochastic with maximum entry \(1/d\).
The character \(\e(kx/n)\) has eigenvalue
\[
 \lambda_k(A_{X_n})
   =\sum_{\substack{a\bmod n\\(a,n)=1}}\e(ka/n)
   =\varphi(n)\frac{\mu(t)}{\varphi(t)},
       \qquad t=n/(n,k).
\]
This is the standard Ramanujan-sum spectrum; we retain its sign.
The only \(t\) with \(\varphi(t)=1\) are \(1,2\).
The eigenvalue at \(t=1\) is the simple constant eigenvalue.
When \(n\) is even, \(t=2\) supplies the unique parity
eigenvalue \(-d\).
Every remaining eigenvalue has modulus at most \(d/2\).
For odd \(n\) there is no \(t=2\), so
\[
 \|P-J_n/n\|\leq1/2.
\]
For even \(n\), all units are odd and the parity parts are equal;
the same bound holds on the complement of the constant
and parity vectors.
These are exact statements for each modulus, not estimates
uniformly transferred from a fixed progression.

We also verify diffuseness without invoking an ineffective
totient estimate.
For a prime-power factor \(p^a\), its factor in
\(\varphi(n)^2/n\) is \(p^{a-2}(p-1)^2\).
It is at least 1 for \(p\geq3\), and the sole possible factor
below 1 is \(1/2\) at \(p=2,a=1\).
Thus \(\varphi(n)\geq\sqrt{n/2}\).
In particular \(1/d\) satisfies every fixed inverse-logarithmic
diffuseness threshold for sufficiently large \(n\).

Use Theorem~\ref{thm:generalweightedentropy} for odd \(n\), and
Theorems~\ref{thm:weightedentropy} and
\ref{thm:bipartitepathentropy} for even \(n\), with a fixed
gap such as \(\delta=1/3\).
Every Hamilton cycle has weight \(d^{-n}\).
Every spanning path has weight \(d^{-(n-1)}\).
The cycle assertion follows directly.
For paths the difference of one \(\log d=O(\log n)\) is
absorbed by the displayed error, uniformly over endpoints.
The totient lower bound makes the error \(o(n)\).
All constants may be chosen absolute because the spectral
gap was uniform over every modulus.
\end{proof}

\paragraph{Maintained spectral input and count conclusions.}
The exact coprime-difference graph, without a prime-modulus or
squarefree-modulus restriction, is formalized in
\texttt{HamiltonianCoprimeFourier}.
Its normalized kernel is nonnegative, symmetric and doubly
stochastic, with maximum entry at most \(1/\varphi(n)\).
\texttt{HamiltonianCoprimeConductor} proves the exact signed
unit-group reduction before any estimate is applied.
\texttt{HamiltonianCoprimeSpectrum} then proves, for every
residue frequency \(k\), the normalized eigenvalue
\[
 \frac{c_n(k)}{\varphi(n)}
    =\frac{\mu(t)}{\varphi(t)},\qquad
 t=\frac n{\gcd(n,k)},
\]
and exactly \(\varphi(t)\) frequencies of each conductor
\(t\mid n\). Conductor one is exactly the constant mode;
conductor two has normalized value \(-1\) and can occur only
at even modulus. Conductors at most two are exactly the
frequencies annihilated by multiplication by two.
Every remaining mode has modulus at most \(1/2\).
\texttt{HamiltonianCoprimeSymbol} proves that the full signed
symbol inverts to the actual physical adjacency kernel.
Using the previously proved real-to-complex norm transfer,
\texttt{HamiltonianCoprimeGap} proves the actual centered
real operator bound \(\|P-J_n/n\|\leq1/2\) for every odd
modulus. Thus the odd spectral hypothesis is discharged,
rather than supplied as an additional assumption.
\texttt{HamiltonianCoprimeEvenGap} also proves the native even
physical half-gap after removing precisely the constant and
parity modes. It identifies the unique conductor-two frequency
with \(n/2\) and its exact inverse-transform parity kernel.
No nonsquarefree modulus or high mode is omitted.
\texttt{HamiltonianCoprimeCrossBlock} now identifies the native
even kernel with its exact bipartite lift and proves stochasticity,
the entry cap and the physical centered half-gap on that cross block.
\texttt{HamiltonianCycleTransport} preserves the full single-cycle
sum under the same vertex equivalence.
\texttt{HamiltonianCoprimeEvenCount} now applies the full
bipartite counting theorem and restores actual arithmetic
weights and reversal classes, uniformly over every even modulus.
The full uniform totient density estimate is compiled in
\texttt{HamiltonianTotientDiffusion}.
\texttt{HamiltonianTotientSublog} supplies its quantified
squared-sublogarithmic consequence.
\texttt{HamiltonianCoprimeAperiodicCount} now applies the
actual odd-modulus spectral gap and all counting hypotheses:
the positive determinant reference and relative error
\(17n^{-1/10}\) hold at every sufficiently large odd modulus.
The full statement and retained normalization are recorded in
Section~\ref{sec:relativecoprime}.
\texttt{HamiltonianCoprimeDeterminant} also proves the exact
full characteristic determinant as the conductor-divisor product
recorded in Section~\ref{sec:relativecoprime}.
Its proof preserves the principal compensation and every
signed mode, including the possible even parity mode.
The positive square-root identification with the explicit
arithmetic constant is now compiled in
\texttt{HamiltonianCoprimeConstant}.
\texttt{HamiltonianCoprimeUnorientedCount} also proves the
actual odd unoriented factorial count, with relative error
at most \(17n^{-1/10}\) uniformly above one threshold.
The full signed logarithmic series and its exact leading
coefficient \(1\) are now compiled in
\texttt{HamiltonianCoprimeConstantSeries}.
\texttt{HamiltonianCoprimeEulerCoefficients} now identifies
every higher coefficient with its exact prime Euler product,
and the full prime-product logarithmic series is compiled.
\texttt{HamiltonianCoprimeConstantBounds} and
\texttt{HamiltonianCoprimeConstantTail} now prove the direct
distinct-prime dependence, uniform geometric remainder and
absolute bounds \(e\leq\mathfrak C(n)\leq e^{11/9}\) over
every odd modulus greater than one.
\texttt{HamiltonianCoprimeConstantLimit} now proves the actual
constant's monotonicity under prime-support inclusion, genuine
multipliability over all odd primes, and the complete primorial
limit with its absolutely convergent infinite-prime series.
Proposition~\ref{prop:coprimeconstant} is fully formalized.
\texttt{HamiltonianCoprimeStirlingCount} now proves the final
odd normalization with its positive displayed reference and
explicit relative error \(23n^{-1/10}\), keeping both the
factorial-count and Stirling errors.
The actual even count is also compiled.
\texttt{HamiltonianCoprimeRelativeCount} combines both parity
branches under one absolute threshold and proves the full
relative arithmetic enumeration with the actual displayed
references and common error at most \(23n^{-1/10}\).
Theorem~\ref{thm:relativecoprime} is fully formalized.
The actual cycle and prescribed-endpoint path entropy conclusions
are now both compiled as follows.

\paragraph{Maintained actual cycle entropy and normalized limit.}
\texttt{HamiltonianCoprimeLogReference} retains the full
parity-dependent logarithmic correction:
\[
 \log D_n-n(\log\varphi(n)-1)
 =\begin{cases}
 \log\pi,&n\ {\rm even},\\
 \log\mathfrak C(n)+
       \tfrac12(\log\pi-\log2-\log n),&n\ {\rm odd}.
 \end{cases}
\]
Its uniform constant bound therefore controls this entire
correction by \(O(\log(2n))\).

\texttt{HamiltonianCoprimeCycleEntropy} consumes the completed
actual all-modulus relative theorem. Above one absolute
threshold it proves \(h(X_n)>0\) and the stronger estimate
\[
 \left|\log h(X_n)-n(\log\varphi(n)-1)\right|
                 \leq K\log(2n),\qquad
 K=4+|\log\pi|+|\log2|.
\]
The relative ratio is eventually within \(1/2\) of one,
so its logarithm contributes at most one; no logarithm of
an unproved positive count is used.
The same owner proves the paper's cycle-error comparison,
using \(n\varphi(n)^{-1/8}\geq1\), and proves the exact
normalized limit formulation
\[
 \forall\varepsilon>0\ \exists N_\varepsilon\
 \forall n\geq N_\varepsilon:\quad
 \left|\frac{\log h(X_n)}n-(\log\varphi(n)-1)\right|
                                      \leq\varepsilon.
\]
All thirteen new declarations pass their focused foundational
audit. Both the cycle assertion and its stated entropy
consequence in Theorem~\ref{thm:coprimecircles} are complete.
\paragraph{Maintained actual integer path entropy.}
The module \texttt{HamiltonianPathCountNormalization} defines the
finite set of prescribed one-orbit path permutations. It proves
its exact membership conditions and identifies its cardinality
\(H_n(a,b)\) with the path sum for the zero--one adjacency.
Its scalar-weight theorem retains precisely \(n-1\) path edges.
The module \texttt{HamiltonianCoprimePathUpper} therefore proves
\[
 Z_{\rm path}(P_n;a,b)=\frac{H_n(a,b)}{\varphi(n)^{n-1}},
 \qquad
 Z_{\rm path}(P_n;a,b)
 \leq\exp\!\left(-(n-1)+16\frac n{\varphi(n)}\log(2n)\right).
\]
The lower count in Section~\ref{sec:relativecoprime} is positive
before any logarithm is taken.
The module \texttt{HamiltonianCoprimePathEntropy} combines the
complete signed lower count with this actual upper bound and
normalization. For one absolute \(N_0\geq512\), every \(n\geq N_0\)
and every compatible distinct pair satisfy \(H_n(a,b)>0\) and
the stronger compiled estimate
\[
 \left|\log H_n(a,b)-n(\log\varphi(n)-1)\right|
     \leq200\,\frac n{\varphi(n)}\log(2n).
\]
The same owner proves the stated weaker bound
\(200n\varphi(n)^{-1/8}\log(2n)\).
It also proves that, for even \(n\), oddness of
\((-(a-b))\bmod n\) is equivalent to opposite parity of the
actual residue representatives. Thus the final theorem has
exactly the prescribed endpoint domain in
Theorem~\ref{thm:coprimecircles}, with no remaining spectral,
count-normalization or endpoint hypothesis.
All ten new membership, count, edge-scaling, upper-count and
entropy declarations pass their combined focused foundational
axiom audit. Every clause of Theorem~\ref{thm:coprimecircles}
is now formalized; the other general weighted, linkage and
layered headline statements retain their own obligations.

\begin{remark}[Normalization, prior existence results, and scope]
When \(n\) is prime, \(X_n\) is complete and
\(h(X_n)=(n-1)!/2\).
When \(n=2^a\), it is complete balanced bipartite, and
\(h(X_n)=(n/2)!(n/2-1)!/2\).
Both exact formulas agree with the normalization above;
cyclic rotations and reversal have been identified, while
translations of the vertex labels have not.

Existence of a Hamilton cycle is not a new conclusion for
this family. Each unit step already gives a Hamilton cycle,
and its opposite gives the same one.
Nor do we claim novelty for the character spectrum or
Hamilton-connectedness.
The quantitative count is supplied to demonstrate a second
arithmetic consumer of the general entropy theorem, now with
the modulus itself growing and with both parity cases covered.
Priority for this counting conclusion requires comparison
beyond the existence and spectral results cited here.
\end{remark}
